\chapter{Functional equations}\index{functional equation}

A functional equation\index{functional equation!definition} asks for
every function satisfying a relation like $f(x+y)=f(x)+f(y)$ --- not a
value, a \emph{family}. Solving one is a two-directions discipline: find
candidates, then prove nothing else works. This chapter is the toolkit:
candidate guessing, substitution tactics, injectivity/surjectivity
milestones, Cauchy's equation done honestly, and the feedback step most
solutions skip --- verifying the answer. After it you can carry a
functional equation from first plug-in to a verified closed form.

\section{What ``solve over $K$'' means}\index{solve over}
``Solve $f(x+y)=f(x)+f(y)$ over $\mathbb Q$'' means: determine
\emph{all} $f:\mathbb Q\to\mathbb Q$ satisfying it, and verify each
claimed solution satisfies it. The domain matters twice --- it decides
which candidates exist (over $\mathbb R$, Cauchy's equation has
infinitely many wild solutions; over $\mathbb Q$, only $kx$) and it
licenses which inductions you may run ($\mathbb N\to\mathbb Z\to\mathbb Q$).

\begin{figure}[ht]\centering
\begin{tikzpicture}[x=1cm,y=0.8cm]
  % nested domains with what f(x+y)=f(x)+f(y) allows on each
  \draw[fill=accent!12,rounded corners=4pt] (0,0) rectangle (10.5,5.4);
  \node[anchor=north west,font=\small\bfseries] at (0.2,5.4) {$\mathbb R$: $f=cx$ \emph{or} pathology (needs a Hamel basis)};
  \draw[fill=accent!25,rounded corners=4pt] (0.4,0.5) rectangle (6.6,3.9);
  \node[anchor=north west,font=\small\bfseries] at (0.6,3.9) {$\mathbb Q$: $f(x)=cx$ forced};
  \draw[fill=accent!45,rounded corners=4pt] (0.8,0.9) rectangle (4.2,2.7);
  \node[anchor=north west,font=\small\bfseries] at (1.0,2.7) {$\mathbb Z$: $f(n)=cn$};
  \draw[fill=accent!65,rounded corners=4pt] (1.2,1.2) rectangle (2.8,2.0);
  \node[font=\small\bfseries] at (2.0,1.6) {$\mathbb N$};
  \node[font=\small,align=right] at (9.9,1.0) {the same equation\\has different answer sets\\on different domains};
\end{tikzpicture}

\caption{The domain ladder for $f(x+y)=f(x)+f(y)$. Notice the answer
set only stabilizes into ``$f=cx$'' where the inductions reach --- on
$\mathbb R$ the tame answer is joined by unclassifiable pathological
solutions.}
\label{fig:fedom}
\end{figure}

\section{Guessing candidates}\index{candidate}
Before proving uniqueness, guess: try $f=0$, $f=c$, $f=kx$, $f=x+c$,
$f=kx+c$, monomials, low-degree polynomials. Each guess that survives is
a lead on the general shape --- and a sanity check for later (if your
``proof'' kills a function you already verified, the proof is wrong).

\section{Substitution tactics}\index{substitution}
The engine of every solution: choose substitutions that make terms
\emph{collide}\index{collide terms}.
\begin{itemize}
\item \emph{Zeros} ($x=0$ or $y=0$) collapse sums and expose constants.
\item \emph{Swaps} ($x\leftrightarrow y$) exploit symmetric left sides to
equate right sides.
\item \emph{Plugs} ($x=u$ for a specific $u$, or $x=f(u)$) force a value
to appear where you need it.
\end{itemize}
Figure~\ref{fig:feflow} is the loop: every substitution should earn a
fact; a substitution that produces $0=0$ taught you nothing --- try a
different one. Figure~\ref{fig:subst} lists the five workhorse
substitutions in the order to try them.
\begin{figure}[ht]\centering
\begin{tikzpicture}[x=1cm,y=1cm,>=stealth]
  \foreach \y/\t in {2.4/{$x=y=0$: constants first},
                    1.4/{$y=0$: isolate $f(x)$ vs $f(0)$},
                    0.4/{$y=-x$: odd/even, cancellation},
                    -0.6/{$x=y$: diagonal $f(2x)$, $f(x)^2$},
                    -1.6/{$y=f(x)$: injectivity probes}}{
    \node[font=\small,anchor=west,fill=graybg,rounded corners=2pt,inner sep=3pt] at (-5.2,\y) {\t};
    \draw[->] (-0.6,\y) -- ({1.4},{0.15*\y});
  }
  \node[draw,rounded corners=3pt,fill=accent!15,font=\small,minimum height=1.6cm,text width=3.0cm,align=center] (eq) at (3.0,0)
    {the equation\\$P(x,y)$};
  \node[draw,rounded corners=3pt,fill=good!20,font=\small,text width=3.2cm,minimum height=1.4cm,align=center] (facts) at (7.4,0)
    {extracted facts:\\$f(0)$, parity, one-variable equation};
  \draw[->,thick] (4.6,0) -- (5.8,0);
\end{tikzpicture}

\caption{The substitution toolbox: each plug is chosen to make two
terms collide and cancel. Notice they are ordered by cheapness ---
zeros first, injectivity probes last.}
\label{fig:subst}
\end{figure}
\begin{figure}[ht]\centering
% Figure: substitution-tactics decision flowchart.
% Source for Figure~\ref{fig:feflow}. TikZ only; regenerable by compiling main.tex.
\begin{tikzpicture}[
  bx/.style={draw=ink,rounded corners=2pt,inner sep=4pt,font=\footnotesize\sffamily,text=slate,align=center},
  lab/.style={font=\footnotesize\sffamily,text=slate,align=center},
  arr/.style={->,draw=rulegray,thick},
]
\node[bx] (q) at (0,0) {read the equation:\\what must collide?};
\node[bx] (z) at (-2.4,-1.3) {plug zeros:\\$x=0$, $y=0$};
\node[bx] (s) at (0,-1.3) {symmetric swap:\\$x\leftrightarrow y$};
\node[bx] (v) at (2.4,-1.3) {plug $x=u$ or $x=f(u)$:\\force a value};
\node[bx] (c) at (0,-2.7) {new fact? loop back ---\\else try an involution $f(f(x))$};
\draw[arr] (q) -- (z);
\draw[arr] (q) -- (s);
\draw[arr] (q) -- (v);
\draw[arr] (z) -- (c);
\draw[arr] (s) -- (c);
\draw[arr] (v) -- (c);
\node[lab] at (0,-3.6) {each substitution should \emph{earn} a fact --- tautologies are waste};
\end{tikzpicture}

\caption{Substitution tactics as a decision loop. Notice the goal of
each move: make two previously-distinct terms coincide.}
\label{fig:feflow}
\end{figure}

\section{Milestones: injectivity, surjectivity, involutions}\index{involution}
Structural facts convert equation-juggling into set-juggling:
$f$ injective lets you cancel $f(\cdot)=f(\cdot)$; surjective lets you
rename any value as $f(\text{something})$; and an
\emph{involution}\index{involution} --- $f(f(x))=x$ --- hands you both
for free (Figure~\ref{fig:invol}).
\begin{figure}[ht]\centering
% Figure: involution-as-bijection arrow diagram.
% Source for Figure~\ref{fig:invol}. TikZ only; regenerable by compiling main.tex.
\begin{tikzpicture}[
  pt/.style={circle,draw=ink,inner sep=1.8pt},
  lab/.style={font=\footnotesize\sffamily,text=slate,align=center},
  arr/.style={<->,draw=accent,thick},
]
\node[pt,label=above:$a$] (a) at (0,1) {};
\node[pt,label=above:$b$] (b) at (1.4,1) {};
\node[pt,label=above:$c$] (c) at (2.8,1) {};
\node[pt,label=above:$d$] (d) at (4.2,1) {};
\draw[arr] (a) to[bend left=25] (b);
\draw[arr] (c) to[bend left=25] (d);
\node[lab,align=center]{$f(f(x))=x$: pairs everything ---\\
injective and surjective for free};
\end{tikzpicture}

\caption{An involution pairs the domain. Notice why it is bijective:
$f(a)=f(b) \Rightarrow f(f(a))=f(f(b)) \Rightarrow a=b$.}
\label{fig:invol}
\end{figure}

\begin{example}[An involution equation, Kyrgyzstan-2012 style]\index{involution!example}
Solve over $\mathbb R$: $f\bigl(x\,f(y)\bigr) = y\,x$.

\textbf{Strategy.} The left side buries $f(y)$ inside $f$ --- expose it.
Plugging $x=1$ should peel out $f(f(y))$; if that comes out as an
involution, bijectivity is free and $y=f(t)$ substitutions become
available.

\textbf{Solution.} $x=1$: $f(f(y)) = y$ --- an involution, so $f$ is a
bijection (Figure~\ref{fig:invol}). Now substitute $y=f(t)$: using
$f\bigl(x\,f(f(t))\bigr) = f(t)\,x$, i.e.\ $f(xt) = x\,f(t)$. Setting
$t=1$: $f(x) = f(1)\,x$, so $f$ is linear, $f(x)=cx$. Feed back:
$c^2\,xy = xy$ for all $x,y$, so $c = \pm 1$. Verification: $f(x)=x$
gives $f(xf(y)) = xy$; $f(x)=-x$ gives $f(-xy) = xy$. \emph{Both} are
solutions --- the answer is the two-element set $\{x\mapsto x,
x\mapsto -x\}$.
\end{example}

\section{Cauchy's equation over $\mathbb Q$}\index{Cauchy equation}
\begin{example}[Solve $f(x+y)=f(x)+f(y)$ over $\mathbb Q$]
\textbf{Strategy.} The candidate is $f=kx$. Rather than prove $f=kx$
directly, expose the free parameter $k=f(1)$ and expand the domain in
steps: $\mathbb N$ by induction, $\mathbb Z$ by a sign trick, $\mathbb Q$ by clearing denominators
(Figure~\ref{fig:stairs}).

\textbf{Solution.} $x=y=0$: $f(0)=0$. Induction gives $f(n)=nf(1)$ on
$\mathbb N$. Then $0=f(n+(-n))=f(n)+f(-n)$ gives $f(-n)=-nf(1)$, so $\mathbb Z$ is done.
For $p/q\in\mathbb Q$: $f(p) = f(q\cdot\tfrac pq) = q\,f(\tfrac pq)$,
so $f(\tfrac pq) = \tfrac pq f(1)$. Hence $f(x)=f(1)x$ for all
$x\in\mathbb Q$ --- and $kx$ does satisfy the equation, so this is the
complete list. Over $\mathbb R$ the pipeline stops: the rationals are
dense but algebra alone doesn't bridge (wild solutions exist without
regularity).
\end{example}
\begin{figure}[ht]\centering
% Figure: domain-expansion staircase N -> Z -> Q -> R for Cauchy.
% Source for Figure~\ref{fig:stairs}. TikZ only; regenerable by compiling main.tex.
\begin{tikzpicture}[
  st/.style={draw=ink,minimum width=1.45cm,minimum height=0.55cm,font=\footnotesize\sffamily,text=slate},
  lab/.style={font=\footnotesize\sffamily,text=slate,align=center},
]
\node[st] (n) at (0,0) {$\mathbb N$: $f(n)=nf(1)$};
\node[st] (z) at (1.7,-0.6) {$\mathbb Z$: $f(-n)=-f(n)$};
\node[st] (q) at (3.4,-1.2) {$\mathbb Q$: $f(p/q)=\frac pq f(1)$};
\node[st] (r) at (5.0,-1.8) {$\mathbb R$: needs a gate*};
\draw[rulegray,thick] (0.85,-0.28) -- (1.05,-0.28) -- (1.05,-0.32);
\draw[rulegray,thick] (2.5,-0.88) -- (2.5,-0.92) -- (2.5,-0.92);
\draw[rulegray,thick] (4.25,-1.48) -- (4.25,-1.52) -- (4.25,-1.52);
\node[lab] at (2.5,-2.6) {*over $\mathbb R$, Cauchy admits wild solutions\\
without an extra hypothesis (monotonicity, boundedness, continuity)};
\end{tikzpicture}

\caption{The domain-expansion staircase for Cauchy's equation. Notice
each step uses the previous one --- and the last step to $\mathbb R$
needs a gate (monotone, bounded, continuous), not more algebra.}
\label{fig:stairs}
\end{figure}

\begin{example}[Jensen's equation by shifting]\index{Jensen}
Solve over $\mathbb Q$: $f(\tfrac{x+y}{2}) = \tfrac{f(x)+f(y)}{2}$.

\textbf{Strategy.} Constants satisfy it --- kill that family first by
shifting: $g(x) = f(x) - f(0)$ satisfies the same equation and has
$g(0)=0$.

\textbf{Solution.} With $g(0)=0$, set $y=0$:
$g(x/2)=g(x)/2$. Then $g(\frac{x+y}{2}) = \frac{g(x)+g(y)}{2}$ becomes
$g(x+y) = g(x)+g(y)$ after replacing $x+y\mapsto$ double --- Cauchy
again, so $g(x)=kx$ and $f(x) = kx + f(0)$. Verify: affine functions
satisfy Jensen. The shift isolated the constant family so the equation
became one already solved.
\end{example}

\section{The pointwise trap}\index{pointwise trap}
Suppose you prove ``for each $t$, $f(t)=\pm t$''. You have \emph{not}
proved $f(t)=t$ or $f(t)=-t$: the sign may depend on $t$
(Figure~\ref{fig:trap}). Kill the gap by mixing: plug $f(a)=a$,
$f(b)=-b$ into the equation and derive a contradiction.
\begin{figure}[ht]\centering
% Figure: pointwise trap -- a piecewise ±x counterexample.
% Source for Figure~\ref{fig:trap}. TikZ only; regenerable by compiling main.tex.
\begin{tikzpicture}[
  lab/.style={font=\footnotesize\sffamily,text=slate,align=center},
  dot/.style={circle,fill=accent,inner sep=1.6pt},
]
\draw[rulegray,->] (-2.4,0) -- (2.4,0);
\draw[rulegray,->] (0,-2.2) -- (0,2.2);
% f(x)=x on x<=0 part drawn solid right-half, f(x)=-x on x>0:
\draw[ink,thick] (-2,2) -- (0,0);
\draw[ink,thick] (0,0) -- (2,-2);
\draw[accent,dashed,thick] (-2,-2) -- (0,0);
\draw[accent,dashed,thick] (0,0) -- (2,2);
\node[lab] at (2.15,1.9) {$+x$};
\node[lab] at (2.15,-1.9) {$-x$};
\node[lab,align=center]{each point obeys $f(t)=\pm t$ --- the sign\\
may change by point, so the global claim fails};
\end{tikzpicture}

\caption{The pointwise trap: $f$ may follow $+x$ on part of the domain
and $-x$ on the rest. Notice this still satisfies $|f(t)|=|t|$ at every
point --- ``per-point'' facts never stitch into a global one.}
\label{fig:trap}
\end{figure}

\begin{example}[Trap autopsy]\index{trap autopsy}
\textbf{Strategy.} In $f(f(x))=x$ + Cauchy-shaped problems the sign
pattern survives mixing only under extra structure. Exhibit the gap:
$f(t) = t$ on $[0,\infty)$, $f(t)=-t$ on $(-\infty,0)$ satisfies
$|f|=|x|$ pointwise --- and fails most equations.

\textbf{Solution.} Take $f(x f(y)) = yx$ again: after proving
$f(f(y))=y$ and $f(x)=f(1)x$ we found $f(1)=\pm1$ --- here the
\emph{parameter} is global, not per-point, so no trap appears. The trap
bites when your deduction is per-$t$ ($f(t)\in\{t,-t\}$) instead of a
single global constant. The check that matters: did the derivation force a
single sign, or leave one free per point?
\end{example}

\section{Feeding back: verify the answer}\index{verification}
``Find all'' ends only when the found functions are plugged back in ---
the verification half that contest graders insist on (the English
handout's rule: never claim without the construction). An equation can
shrink a candidate family to $\{f=0\}$, and even that last member still
has to be checked.

\begin{example}[Verification eliminates a wrong guess]
Suppose solving $f(xy)=f(x)f(y)$ over $\mathbb R_{>0}$ led you to the
guess $f(x)=x+c$. Show verification rejects every $c\neq0$.
\par\smallskip
\textbf{Strategy.} Plug the candidate in literally --- don't argue
about whether it ``should'' work.
\par\smallskip
\begin{solution}
$f(xy)=xy+c$ while $f(x)f(y)=(x+c)(y+c)=xy+c(x+y)+c^2$. Equality
requires $c(x+y)+c^2=c$ for all $x,y$; fixing $y$ and varying $x$
forces $c=0$ (then $c^2=c$ is automatic). The guess is rejected ---
and verification catches what the solving pass missed: the real answer
family is $f \equiv 0$ together with $f(x)=x^k$ (under any mild
regularity hypothesis).
\end{solution}
\end{example}

\section{Summary}
\begin{center}\small
\begin{tabular}{@{}ll@{}}
\toprule
\textbf{Tactic} & \textbf{What it earns} \\
\midrule
Plug $0$ & constants, $f(0)$ \\
Swap $x\leftrightarrow y$ & equality of two right sides \\
Plug $x=u$, $x=f(u)$ & forced values, involutions \\
Induction & $\mathbb N$ expansion \\
Sign trick $f(-x)$ & $\mathbb Z$ expansion \\
Denominator clear & $\mathbb Q$ expansion \\
Shift $g=f-f(0)$ & kill constant families \\
Mix two points & kill the pointwise trap \\
Plug back in & verification (required half) \\
\bottomrule
\end{tabular}
\end{center}

\section{Code: hunting solutions by brute force}\index{fe\_search}
\texttt{fe\_search.py} scans degree-$\le2$ polynomials
$f(t)=At^2+Bt+C$ on a grid for the equation
$f(x+y)-f(x-y)=4xy$.
\begin{checkpoint}
Predict which coefficient must be $1$, which must be $0$, and which is
free. Then run \texttt{python3 fe\_search.py}.
\end{checkpoint}
\begin{lstlisting}
quadratic survivors: [(1.0, 0.0)]
constant term free? distinct C values per (A,B): True
  f=t^2 + c: residual 0.000
  f=t^2 - 1: residual 0.000
  f=2t^2: residual 2500.000
  f=t: residual 2550.000
  f=0: residual 2500.000
verdict: f(t) = t^2 + C is the whole degree-<=2 family
\end{lstlisting}
The brute force agrees with the algebra ($f(x+y)-f(x-y)=4xy$ is
$f=t^2+C$ exactly) --- and confirms the constant family survives every
equation of difference type. \texttt{test\_14\_fe.sh} re-runs and greps
the verdict line.

\section{Common mistakes}\index{functional equation!mistakes}
\begin{itemize}
\item \emph{Never verifying.} ``Find all'' is two directions; a
candidate list without a plug-back-in check is half a solution.
\item \emph{The pointwise sign.} $f(t)=\pm t$ per point $\ne$ a global
$\pm$. The sign must come out as a single constant.
\item \emph{Tautology loops.} Substitutions that yield $0=0$ produce no
fact --- if three plugs teach nothing, change \emph{what} you plug, not
how hard you stare.
\end{itemize}

\section*{Exercises}
\subsection*{Inline}
\begin{exercise}[easy]\label{exr:f1}
From $f(x+y)=f(x)+f(y)$ deduce $f(0)=0$ and $f(-x)=-f(x)$.
\end{exercise}
\begin{solution}
$x=y=0$ gives $f(0)=2f(0)$. Then $0=f(0)=f(x)+f(-x)$.
\end{solution}
\begin{exercise}[easy]\label{exr:f2}
Show $f(x)=kx$ satisfies $f(x+y)=f(x)+f(y)$ over $\mathbb R$ --- the
verification half.
\end{exercise}
\begin{solution}
$kx+ky = k(x+y)$. Done --- but it is \emph{required}, not optional.
\end{solution}
\begin{exercise}[medium]\label{exr:f3}
Solve $f(x+y)=f(x)f(y)$ over $\mathbb Q$ with $f\not\equiv0$, given $f(1)>0$.
\end{exercise}
\begin{solution}
$f(0)=1$ (since $f(x)=f(x)f(0)$ and some value is nonzero). Induction:
$f(n)=f(1)^n$; $f(-n)=f(1)^{-n}$; $f(p/q)=f(1)^{p/q}$ with the rational
root defined since $f(1)>0$. Answer $f(x)=a^x$ for $a=f(1)>0$ ---
verify: $a^{x+y}=a^xa^y$.
\end{solution}
\begin{exercise}[medium]\label{exr:f4}
Exhibit a function $|f(t)|=|t|$ everywhere that is neither $t$ nor $-t$,
and use it to explain the pointwise trap.
\end{exercise}
\begin{solution}
$f(t)=t$ on $[0,\infty)$, $f(t)=-t$ on $(-\infty,0)$ works (as in
Figure~\ref{fig:trap}); the sign is per-point, so a ``$f(t)\in\{t,-t\}$
for each $t$'' conclusion cannot be globalized without a mixing argument.
\end{solution}
\subsection*{Additional}
\begin{exercise}[easy]\label{exr:fa1}
Solve $f(x+1)=f(x)+2$ over $\mathbb Z$ with $f(0)=3$.
\end{exercise}
\begin{solution}
$f(n)=2n+3$ by induction both directions; verify.
\end{solution}
\begin{exercise}[easy]\label{exr:fa2}
Which constant functions satisfy $f(x+y)=f(x)+f(y)+7$?
\end{exercise}
\begin{solution}
$f=c$: $c=2c+7 \Rightarrow c=-7$. Verify: $-7=-7-7+7$.
\end{solution}
\begin{exercise}[easy]\label{exr:fa3}
Show an involution is bijective.
\end{exercise}
\begin{solution}
$f(f(x))=x$: injective since $f(a)=f(b)\Rightarrow a=f(f(a))=f(f(b))=b$;
surjective since $x=f(f(x))$ is in the image.
\end{solution}
\begin{exercise}[easy]\label{exr:fa4}
Solve $f(f(x))=x$ and $f(x)=kx$ simultaneously over $\mathbb R$.
\end{exercise}
\begin{solution}
$k^2=1$, so $f=x$ or $f=-x$; both check.
\end{solution}
\begin{exercise}[medium]\label{exr:fa5}
Solve over $\mathbb Q$: $f(x+y)=f(x)+f(y)+xy$.
\end{exercise}
\begin{solution}
$g=f-x^2/2$ is additive on $\mathbb Q$: $g(x)=kx$; $f=kx+x^2/2$. Verify.
\end{solution}
\begin{exercise}[medium]\label{exr:fa6}
Prove $f(f(x))=x$ has no solution of the form $f(x)=x+c$ unless $c=0$.
\end{exercise}
\begin{solution}
$f(f(x))=x+2c=x$ needs $c=0$.
\end{solution}
\begin{exercise}[medium]\label{exr:fa7}
Suppose $f(x+y)=f(x)+f(y)$ over $\mathbb R$ and $f$ is monotone. Prove
$f=kx$ (the gate of Figure~\ref{fig:stairs}).
\end{exercise}
\begin{solution}
On $\mathbb Q$, $f=kx$ with $k=f(1)$. For $r\in\mathbb R$, pick rationals
$q_n\uparrow r$, $q'_n\downarrow r$; monotonicity sandwiches
$f(r)$ between $kq_n, kq'_n \to kr$. Monotonicity then forces
$k = f(1) \ge 0$.
\end{solution}
\begin{exercise}[medium]\label{exr:fa8}
Describe all solutions of $f(xy)=f(x)+f(y)$ over positive rationals
with $f(2)=1$.
\end{exercise}
\begin{solution}
On the multiplicative group $\mathbb Q_{>0}$, $f$ is determined by its
values on primes: $f(x)=\sum_p e_p\,c_p$ for $x=\prod p^{e_p}$, with
$c_2 = f(2) = 1$ and the other $c_p$ free --- so $f=\log_2$ is only
the special case $c_p=0$ for $p>2$, not the general answer.
\end{solution}
\begin{exercise}[hard]\label{exr:fa9}
Prove Cauchy's equation over $\mathbb R$ admits non-linear solutions
(hence the gate) using a Hamel basis.
\end{exercise}
\begin{solution}
Pick a $\mathbb Q$-vector-space basis $H$ of $\mathbb R$ containing $1$.
Define $f$ linear on basis elements with $f(1)=1$, $f(h)=0$ for
$h\ne1$; extend $\mathbb Q$-linearly. Then $f$ is additive but not $kx$
(e.g.\ $f(\sqrt2)=0$).
\end{solution}
\begin{exercise}[hard]\label{exr:fa10}
Implement-or-derive: extend \texttt{fe\_search.py} to degree-3
polynomials on the same equation and report the surviving family.
\end{exercise}
\begin{solution}
The same $t^2+C$ survives; no cubic term can cancel ($A t^3$'s residual
grows like $|t|^3$). Brute force confirms the degree bound.
\end{solution}
\begin{exercise}[easy]\label{exr:fa11}
Solve $f(x+1)=f(x)+2$ over $\mathbb Z$.
\end{exercise}
\begin{solution}
Induction both directions: $f(n)=2n+f(0)$; every such function works.
\end{solution}
\begin{exercise}[medium]\label{exr:fa12}
Find all $f:\mathbb Q\to\mathbb Q$ with $f(x+y)=f(x)f(y)$.
\end{exercise}
\begin{solution}
$f(0)=f(0)^2$, so $f(0)\in\{0,1\}$; $f(0)=0$ forces $f\equiv0$. Else
$f(0)=1$ and $f(x)f(-x)=1$, so $f$ is never $0$, and
$f(x)=f(x/2)^2>0$. With $a=f(1)$: $f(n)=a^n$, and
$f(1)=f(1/q)^q$ forces $a^{1/q}\in\mathbb Q$ for \emph{every}
denominator $q$ --- impossible for $a\neq1$ (write $a=r/s$ reduced;
$v_p(a)$ would need to be divisible by every $q$). So $a=1$: the only
solutions are $f\equiv0$ and $f\equiv1$.
\end{solution}
\begin{exercise}[medium]\label{exr:fa13}
Solve $f\bigl(x+f(y)\bigr)=x+y$ over $\mathbb R$.
\end{exercise}
\begin{solution}
$x=0$ gives $f(f(y))=y$: an involution, hence bijective. $y=0$ gives
$f(x+f(0))=x$, i.e.\ $f(t)=t-c$ with $c=f(0)$. Plug into the
involution: $t-2c=t$ forces $c=0$, so $f(x)=x$; it verifies.
\end{solution}
\begin{exercise}[hard]\label{exr:fa14}
Show no \emph{continuous} $f:\mathbb R\to\mathbb R$ satisfies
$f(f(x))=-x$, then exhibit one over $\mathbb C$.
\end{exercise}
\begin{solution}
$f\circ f=-\mathrm{id}$ is bijective, so $f$ is; a continuous bijection
is strictly monotone, and the composite of two monotone maps (same or
opposite direction) is increasing --- but $x\mapsto -x$ is strictly
decreasing. Contradiction. Over $\mathbb C$: $f(z)=iz$.
\end{solution}
\begin{exercise}[hard]\label{exr:fa15}
Show every solution of $f(x+y)=f(x)+f(y)$ over $\mathbb R$ that is
bounded on some interval is linear --- the regularity bridge from
$\mathbb Q$ to $\mathbb R$.
\end{exercise}
\begin{solution}
$g(x)=f(x)-f(1)x$ is additive and vanishes on $\mathbb Q$; boundedness
of $f$ on an interval bounds $g$ near $0$, and an additive function
bounded near $0$ is continuous at $0$ (since $g(nh)=ng(h)$, unbounded
blow-up otherwise), hence continuous everywhere, hence $g\equiv0$ ---
so $f(x)=f(1)x$.
\end{solution}

\section*{Notes}
The candidates→substitutions→involution→Cauchy backbone follows Chen's
\emph{Introduction to Functional Equations}~\cite{chen-funceq}; the
$\mathbb N\to\mathbb Z\to\mathbb Q$ pipeline and the $\mathbb R$-gate are the Excerpts'
treatment~\cite{otis}; Engel's chapter supplies the Jensen and
involution stock~\cite{engel}. The pointwise-trap framing is the
``English'' handout's two-directions rule~\cite{chen-english} applied to
sign choice. Contest uses: 2023 AMC 12A P22, P25; 2015 AIME I P10;
2010 AIME I P6; 2024 AIME I P12; USAJMO 2024 P5; JMO 2015 P4; USAMO
2002 (read inside~\cite{chen-funceq}); IMO 2026 P5.

\section*{Further reading}
\begin{itemize}
\item E.~Chen, \emph{Introduction to Functional Equations} --- free
PDF; the reference this chapter follows:
\url{https://web.evanchen.cc/handouts/FuncEq-Intro/FuncEq-Intro.pdf}
\item E.~Chen, \emph{The English Handout} --- how to write
``determine all'' answers in both directions:
\url{https://web.evanchen.cc/handouts/english/english.pdf}
\item A.~Engel, \emph{Problem-Solving Strategies}, functional-equations
chapter --- the drill stock.
\item UNSW CPMSoc workshop slides (the functional-equations PDF is
under \emph{Mathematics 2021}):
\url{https://unswcpmsoc.com/workshops/}
\item R.~Hammack, \emph{Book of Proof} --- for the
function/bijection machinery the milestones section reuses:

\url{https://richardhammack.github.io/BookOfProof/}
\item Y.~Zhao's olympiad handouts --- functional-equation problems in
the olympiad style:
\url{https://yufeizhao.com/olympiad/}
\end{itemize}

\begin{rulebox}[Watch instead]
\begin{itemize}
\item \textbf{Michael Penn --- ``A functional equation that didn't
quite make the IMO''}, 6:43,
\url{https://www.youtube.com/watch?v=ux9KYSK0jl8}
--- solves
$f(x+y)+f(x-y)=f(x)f(y)$ by exactly the substitution tactics of
the Substitution-tactics section: creative value choices, extracting
$f(0)$, peeling to one variable. (Verified 2026-10-11.)
\end{itemize}
\end{rulebox}

\chapternotes
